\chapter{数据设计}

\section{数据对象说明}
LingoBridge 平台底层的核心实体对象在后端 Repository 仓库中抽象为强类型的数据传输对象（DTO），确保在双模存储环境下数据结构的完整性。系统运行时的元数据主要涵盖用户账户数据、班级组织数据、班级成员关联数据、课程大纲数据、课时节点数据、存储文件数据、作业任务数据以及学习成果评估记录。这些对象均遵循唯一性约束与级联删除规则，防止产生无主孤立节点。数据对象的字段级物理表定义（如字段约束、数据长度、索引配置等）属于详细设计阶段，本说明书仅从高层逻辑模型上规范数据实体关联。

\section{ER 模型设计}
系统概念结构设计与实体关系（ER）模型完全对齐平台的业务流，通过外键约束维持数据实体间的一致性。LingoBridge 系统核心数据库实体关系图如图 5.1 所示。

\begin{figure}[H]
\centering
\includegraphics[width=0.78\textwidth]{figures/LingoBridge关系图.pdf}
\caption{LingoBridge 系统核心实体关系图（ERD）}
\label{fig:erd_main}
\end{figure}

如图 5.1 所示，系统最顶层实体为用户表（\texttt{users}），在此基础上通过班级表（\texttt{classes}）和班级成员关联表（\texttt{class\_members}）实现师生关系的管理。一个班级可以绑定多个课程（\texttt{courses}），而一个课程又划分为多个课时节点（\texttt{lesson\_nodes}）。每个课时节点在物理上可以关联一个默认的课件文件（\texttt{files}），并包含多个通过 Excel 导入的作业任务项（\texttt{learning\_tasks}）。最终，学生的每次听读作答数据都会沉淀为学习记录（\texttt{learning\_records}），实现评估闭环。

\section{数据表逻辑实体定义}
为了指导物理建表并与详细设计进行职责交接，本节从概念实体维度明确各数据表的核心职责分工。

一是用户身份实体，主要承载学生与教师的登录凭证与多语种基础属性，是系统鉴权与业务多租户绑定的逻辑根节点。

二是班级与教务关联实体，主要维护试点学校小班级组织划分，提供学生加入与剔除班级的持久化入口，以支持教师对学生教务日程的分发。

三是课程与大纲课时实体，主要承载课程结构、发布状态以及具体排课时间段的管理，与前端的课程日历日程组件保持逻辑映射。

四是媒体文件及课件作业包实体，主要维护 PDF 课件元数据以及由 Excel 导入的词汇听写跟读题目白名单。

五是学习评测成果记录实体，主要记录学生每次课后练习的最佳音视频成绩、答题次数以及完成标记，是后续教学质量评估的底层基准数据。

\section{数据流设计}
平台在进行动态业务操作时，数据流严格遵循单向传递原则。当学生上传听读 Opus 音频时，前端子系统通过二进制流将音频封装，随 DTO 上报至后端 Express 的 \texttt{recordings} 路由。后端接收到数据后，首先向文件服务模块申请随机 UUID 文件标识并登记入 \texttt{files} 表中，同时将物理文件写入磁盘缓冲区。随后，后端将文件元数据和学生请求组装，发送至语音服务模块，得到音素级的评估分数。后端经过逻辑计算，比对 \texttt{learning\_records} 中存储的历史最佳成绩，如果本次评测的分数高于历史最佳分，则执行安全无损更新，最终将最新成绩与评测反馈包返回给前端，完成单次音频数据评测的完整闭环。

\section{数据安全与备份设计}
为了确保系统试点元数据和留学生个人档案数据的物理安全，数据库底层配置了多层安全与备份架构。

在物理传输维度，PostgreSQL 开启了标准的强制加密传输协议（SSL/TLS），保障数据库在长途数据传输时的抗监听性能。

在逻辑控制层面，数据仓库模块为多租户配置了严格的行级安全控制（RLS）和外键级联策略，确保学生角色在拉取学习记录时，绝无可能非法查询到非本人或非本班的敏感数据。

在灾备维度，后端部署脚本中内置了 PM2 驱动的定时数据物理导流计划，每日凌晨自动调度容器执行备份，生成带时间戳的安全物理 SQL 快照，流式归档到独立的磁盘缓冲区，确保灾后数据库能够在较短时间内实现无损热重建。
